\section*{5. 附录: 整性定理}

\begin{center}
P. Deligne
\end{center}

\subsection*{5.0.}
设 $p$ 为素数, $q$ 为 $p$ 的幂, $\mathbb F_q$ 为有 $q$ 个元素的域,
$E$ 为数域在素数 $\ell\ne p$ 之上的一个位处的完备化, $T$ 为素数的集合
(例如 $T=\varnothing$).

$E$ 的代数闭包中的元素 $\alpha$ 称为\emph{$T$-整的 (T-integral)},
如果它在 $\mathbb Q$ 上代数, 且在 $\mathbb Z[(1/r)_{r\in T}]$ 上整.
有限维 $E$ 向量空间的自同态 $F$ 称为 $T$-整的, 如果其特征值均为 $T$-整的.

设 $k$ 是有 $q'$ 个元素的有限域, $\bar k$ 是 $k$ 的代数闭包, 设
\[
\mathfrak G=\operatorname{Gal}(\bar k/k)
\]
为 $k$ 的 Galois 群. 用 $f\in\mathfrak G$ 表示 Frobenius 置换
\[
f(x)=x^{q'},
\]
用 $F$ 表示“几何” Frobenius, 即 $f$ 在 $\mathfrak G$ 中的逆.
$\mathfrak G$ 在有限维 $E$ 向量空间上的表示 $\rho$ 称为 $T$-整的,
如果 $\rho(F)$ 是 $T$-整的.

若 $X$ 是 $\mathbb F_q$ 上的有限型概形, $x$ 是 $X$ 的闭点,
则以 $F_x$ 表示 $k(x)$ 的 Galois 群中的几何 Frobenius 元素.
若 $\bar x$ 是 $X$ 中位于 $x$ 之上的几何点, $\underline G$ 是 $X$ 上的 $E$ 层,
则通过结构搬移, $F_x$ 作用于纤维 $G_{\bar x}$.

\subsubsection*{定义 5.1.}
$\mathbb F_q$ 上有限型概形 $X$ 上的可构造 $E$ 层 $\underline G$ 称为 $T$-整的,
如果对 $X$ 的每个闭点 $x$ 之上的 $X$ 的每个几何点 $\bar x$, $F_x$ 都是 $G_{\bar x}$ 的
$T$-整自同构.

若 $a:X\longrightarrow\operatorname{Spec}(\mathbb F_q)$ 是 $\mathbb F_q$ 上的有限型概形,
$\underline G$ 是 $X$ 上的可构造 $E$ 层, 则用 $\underline H^*(X,\underline G)$
表示 $\operatorname{Spec}(\mathbb F_q)_{\acute e t}$ 上的 $E$ 层 $R^*a_*\underline G$.
若 $\overline{\mathbb F}_q$ 为 $\mathbb F_q$ 的代数闭包, 且
$\overline X=X\times_{\mathbb F_q}\overline{\mathbb F}_q$, 则这些层可与作为 Galois 模的
上同调群 $H^*(\overline X,\underline G)$ 等同, 即赋予其几何 Frobenius $F$ 的作用.
对固有支撑上同调采用同样约定.

\subsubsection*{引理 5.2.1.}
设 $X$ 是 $\mathbb F_q$ 上维数至多为 $n$ 的有限型概形,
$\underline G$ 是其上的可构造 $E$ 层.

\begin{enumerate}[label=(\roman*)]
\item 存在 $X$ 的维数为 $0$ 的闭子集 $Y$, 使显然的映射
  \[
  \underline H^0(X,\underline G)\longrightarrow
  \underline H^0(Y,\underline G)
  \]
  为单射.
\item 若 $X$ 分离, $U$ 是 $X$ 的开子集, 其补集 $Z$ 维数小于 $n$,
  则存在 $U$ 中维数为 $0$ 的闭子集 $Y$, 以及满射
  \[
  \underline H^0(Y,\underline G)(-n)\longrightarrow
  \underline H_c^{2n}(X,\underline G).
  \]
\end{enumerate}

\medskip
\noindent\emph{证明.}
断言 (i) 显然. 现证 (ii). 将 $X$ 换成 $X_{\mathrm{red}}$ 并缩小 $U$ 后,
可设 $U$ 光滑且 $\underline G|_U$ 为扭常值层. 由于 $\dim Z<n$, 有
\[
0=\underline H_c^{2n-1}(Z,\underline G)
\longrightarrow \underline H_c^{2n}(U,\underline G)
\xrightarrow{\sim}\underline H_c^{2n}(X,\underline G)
\longrightarrow \underline H_c^{2n}(Z,\underline G)=0,
\]
而 Poincar\'e 对偶 (SGA~4, Expos\'e~XVIII 及 [2]) 给出
\[
\underline H_c^{2n}(U,\underline G)\simeq
\underline H^0(U,\underline G^\vee(n))^\vee.
\]
结论由 (i) 得出.
